\section{A companion erratum to the printed positive form}
\label{sec:positive-form-erratum}

\paragraph{Inherited finding P-F12 / PL-Q22, rechecked.}
In GCT VI, Theorem 1.6.1(d) and equation (4.4), the numerator
index cannot literally be the degree $d$ of the stretching
quasipolynomial, with numerator degree bounded by $d$.
The Kahle--Micha\l ek ray already contradicts that reading.
The correct Cohen--Macaulay positive form has unrestricted
numerator degree; the denominator has $d+1$ factors when
multiplicities are counted.  This correction does not rescue
PH3, whose two clauses fail by the cubic theorem proved above.

For clarity, all data needed for the contradiction are
\[
 \phi(s)=\frac{s+r_s}{3},\qquad
 (r_0,\ldots,r_5)=(3,-1,1,0,2,-2),\qquad
 \sum_{s\ge0}\phi(s)t^s=\frac{1+t^4}{(1-t^2)(1-t^3)}.
 \tag{PNEW-KM}
\]
This is Kahle and Micha\l ek (2018), Theorem 5 and Remark 6,
printed p.~4, for $a_{(4s),(3)}^{(7s,5s)}$.
Here is also a direct coefficient verification.  Put
$d(j)=\#\{(a,b)\in\mathbb Z_{\ge0}^2:2a+3b=j\}$,
with $d(j)=0$ for $j<0$.  Weight counting in
$\operatorname{Sym}^{4s}(\operatorname{Sym}^3\mathbb C^2)$
gives the Gaussian polynomial
\[
 \prod_{i=1}^3\frac{1-q^{4s+i}}{1-q^i}.
\]
Indeed its coefficient counts partitions in the $3$ by $4s$
rectangle, equivalently the monomial weight multiplicities.
For an $\mathfrak{sl}_2$ module, the multiplicity of highest
weight $2s$ is the difference between the weight-$2s$ and
weight-$(2s+2)$ multiplicities, since each irreducible has
one copy of every weight from its highest weight down by twos.
Multiplication of the Gaussian polynomial by $1-q$ therefore
gives
\[
 \phi(s)=d(5s)-d(s-1)-d(s-2)-d(s-3)\quad(s\ge1).
\]
Terms with two numerator shifts have degree greater than $5s$.
The elementary formula
$d(6a+r)=a+1-\mathbf1_{r=1}$ for $a\ge0$, $0\le r<6$,
follows by choosing $b$ of the same parity as the argument.
Substitution proves (PNEW-KM) for $s\ge3$; $s=0,1,2$
are checked directly.  Summing
$\phi(s)=\lfloor s/3\rfloor+\mathbf1_{2\mid s}$ gives
the rational series in (PNEW-KM).

The degree is one.  A positive numerator of degree at most one
and a denominator with two factors would have the form
\[
                \frac{1+ct}{(1-t^a)(1-t^b)},\qquad c\ge0.
\]
Since $\phi(1)=0$, one must have $c=0$ and $a,b>1$.
Then $\phi(2)=\phi(3)=1$ forces $\{a,b\}=\{2,3\}$.
The resulting coefficient of $t^4$ is one, whereas
$\phi(4)=2$, a contradiction.  The unrestricted positive
numerator $1+t^4$ above is consistent with
Cohen--Macaulayness.

For later use, (PNEW-KM) also gives exactly
$p(\phi)=s(\phi)=2$.  Shifting by two makes every constituent
coefficient nonnegative.  Before that shift the residue-five
constituent $(X-2)/3$ is nonpositive at a positive integer:
its values at 1 and 2 are negative and zero, respectively.
This proves both lower bounds and the matching upper bounds.
